\subsection{Overview}

The evaluation of models used to simulate and design agricultural subsurface (tile) drainage rests on several converging lines of evidence: (i) site-scale field validations that compare simulated water-table depths, drain discharge and, increasingly, salinity and solute fluxes against monitored records; (ii) benchmark comparisons of analytical or semi-analytical drain-spacing solutions against laboratory experiments, numerical codes and one another; (iii) watershed- and regional-scale studies that test how drainage is conceptualised within distributed groundwater and integrated hydrological models; and (iv) formal sensitivity, calibration and uncertainty analyses. Across this body of work several robust conclusions emerge. The field-scale water-management model \textsc{DRAINMOD} is by far the most extensively tested and is repeatedly judged a reliable design and evaluation tool for poorly and artificially drained lands \parencite{Singh2019b,Gowda2012}. Model skill is strongly scale-dependent in time: monthly and annual drainage volumes and cumulative water balances are typically reproduced with high accuracy, whereas daily and sub-daily peak flows, low continuous dry-season discharge, and snowmelt events are the persistent weak points \parencite{Gowda2012,Dayyani2009,Morrison2014,Varvaris2018}. Saturated flow-based formulations can substantially over-predict drainage relative to vadose-zone-aware analyses, sometimes by an order of magnitude, so capillary and unsaturated-zone storage cannot be neglected \parencite{Xin2016b}. Saturated hydraulic conductivity and the method used to estimate evapotranspiration dominate model sensitivity and prediction uncertainty \parencite{NabiSichani2012,Malakshahi2020,Wang2006}. At the catchment scale, the way tile drains and subsurface heterogeneity are represented can shift the simulated drain contribution to streamflow by a factor of two or more \parencite{Karan2021,Boico2022}. This section synthesises the quantified performance, mechanisms and design implications established by these studies.

\subsection{DRAINMOD as the Benchmark Field-Scale Model}

\textsc{DRAINMOD} is described as one of the first widely adopted and most extensively used models for assessing drainage performance and simulating subsurface water management in tile-drained fields \parencite{Singh2019b}. A recurring, well-supported finding is that its process-based structure, with independently measurable inputs, allows good description of drainage hydrology across a wide range of climates when soil hydraulic properties are adequately characterised \parencite{Gowda2012}.

\subsubsection{Water Table, Drain Discharge and Temporal Scale}

Across humid and temperate sites, \textsc{DRAINMOD} reproduces water-table depth and drainage discharge with generally close agreement, but its accuracy improves markedly from the daily to the monthly and annual scale. In North Carolina, validation Nash--Sutcliffe efficiencies (EF) for water-table depth were 0.68--0.72, daily drainage EF fell to 0.49--0.73, monthly drainage EF rose to 0.77--0.87, and predicted annual drainage volumes lay within 3 cm (5\%) of observations with a two-year cumulative difference of just 1 cm (1\%) \parencite{Gowda2012}. Similar behaviour was found in Iowa, where predicted and observed monthly subsurface drainage matched closely and cumulative 1990--2003 drainage was within 2\% (WEBS\_CC) to 10\% (CANI\_CS) of measured values, with a slight tendency to over-predict in low-drainage years \parencite{Singh2006b}. Two-year evaluations in eastern Canada reported simulated annual drainage about 30 mm below observed while capturing timing and magnitude of events \parencite{Dayyani2009}. Repeated across studies is the tendency to underestimate daily peak flows even where general trends are well captured \parencite{Morrison2014,Singh1996}. Uncertainty analysis in Indiana confirmed this scale dependence: 90\% confidence-interval widths for drain flow grew from 0--0.6 cm day$^{-1}$ (daily) and 0--3.1 cm month$^{-1}$ (monthly) to 7.6--12.4 cm year$^{-1}$ (yearly), yet annual DRAINMOD predictions fell well within the bounds \parencite{Wang2006}.

\subsubsection{Performance in Cold and Semi-Arid Climates}

\textsc{DRAINMOD} has been validated well beyond its original humid setting. In cold regions of eastern Canada it accurately simulated the timing and magnitude of subsurface drainage and water-table fluctuations, including winter months and spring snowmelt, with calibration/validation $R^2$ up to 0.93 for water-table depth \parencite{Dayyani2009,Morrison2014}. Calibration substantially improved monthly performance, although surface runoff was predicted more poorly than subsurface drainage and daily peaks were underestimated, with better performance where soils were uniform than where cracking and preferential flow occurred \parencite{Morrison2014}. Under semi-arid Indian conditions, measured and predicted midspan water-table heights generally agreed (average absolute deviations of roughly 14--24 cm), though simulated water tables responded faster to rainfall and the model over-estimated depths in dry periods and under-estimated them in wet periods; the authors nonetheless judged the model satisfactory for drainage design in such climates \parencite{GUPTA1993}. On Illinois claypan soils, measured and simulated weekly cumulative SEW30 agreed with a coefficient of determination of 0.91, and subsurface drainage was shown to affect trafficability, dry days and SEW30 far more than surface drainage \parencite{MOSTAGHIMI1989}. Additional acceptable-to-good calibrations (e.g. $R^2=0.89$--0.93, NSE $=0.84$--0.89) have been reported for groundwater-table prediction in further settings \parencite{Abduljaleel2023,Awad2022b}.

\subsubsection{Documented Limitations and Failure Modes}

The literature also delineates conditions under which \textsc{DRAINMOD} fails. At an Israeli site it showed very large deviations, with water-table differences of 0.3--1.7 m in four of five replicates and reasonable agreement in only one plot; a detailed piezometric study attributed the disagreement to a river boundary condition, permeable sandy layers below drain depth, and deep plus lateral seepage exceeding 50\% of annual irrigation plus rainfall \parencite{Sinai2006}. A structural drawback is the inability to represent time-variable surface depression storage: with a fixed depression-capacity input the model assumed excess water infiltrated, producing an unrealistic water fate and a marked observed--predicted gap during a mid-summer paddy drainage period, which may compromise nutrient-fate predictions \parencite{Abduljaleel2023}. The model tends to over-predict drainage during low-drainage years \parencite{Singh2006b} and, being built around parallel-drain geometry, requires plot-by-plot calibration when adapted to landscapes without a regular drain network, where equivalent drain-spacing parameters ranged from about 1700 to 4000 cm and porous organic layers strongly influenced water tables \parencite{He2002}.

\subsubsection{Parameterisation, Pedotransfer Functions and Enhancements}

Because site-specific hydraulic data are often lacking, several studies established that pedotransfer estimation combined with automated calibration can drive \textsc{DRAINMOD} reliably. Coupling ROSETTA-derived properties with the PEST estimator produced calibrated, validated Iowa applications \parencite{Singh2006b}, and three levels of soil hydraulic parameter sets (including ROSETTA with a soil survey database) all gave acceptable annual and monthly drainage, with only subtle differences among them, although calibrated parameters that were best in calibration performed worst in validation \parencite{Qi2015}. In South Africa, van Genuchten fits to retention data ($R^2>0.90$) supported validation NSE of 0.86 (water-table depth) and 0.81 (discharge) \parencite{Malota2024,Malota2015}. Enhancements have targeted the daily-peak weakness: incorporating macropore flow into \textsc{DRAINMOD} improved daily drainage-peak prediction, captured about 75\% of 60 peak events, partitioned matrix and macropore flow (macropore flow on 98 days), and predicted cumulative drainage within 8\% (65.9 vs 71.5 cm), though surface runoff was strongly over-predicted (51.9 vs 23.7 cm) and annual budgets closed only to about 6.4 cm \parencite{Askar2020}. A neural-network surrogate (\textsc{DRAINMOD-ANN}) successfully cloned daily groundwater-table predictions across drainage layouts, with GRNN outperforming multilayer feed-forward training, offering a way to enhance performance without altering source code \parencite{Awad2022,Awad2022b}. Sensitivity to the evapotranspiration algorithm is pronounced: replacing the default Thornthwaite method with Rohwer or pan-evaporation-based estimates substantially improved water-table prediction in northern Iranian paddies, and the best ET-estimation methods were not necessarily the best for water-table simulation \parencite{Malakshahi2020}. Likewise, moving ET$_0$ computation from a daily to an hourly time step shifted the simulated water fate, increasing cumulative runoff by 4.8\% and decreasing drainage and infiltration by 3.1\% and 1\% \parencite{Awad2022b}.

\subsection{Other Field-Scale Physically Based Models}

\subsubsection{Richards-Equation and Variably Saturated Models}

The HYDRUS (2D/3D) code is consistently found capable of predicting drainage fluxes and water-table depths under contrasting subsurface systems, with strong calibration and validation skill (EF from 0.91 to 1; RMSE 0.2--6.4 cm) in Iranian paddy fields, where it is recommended for identifying optimal drainage layouts \parencite{Karandish2018,DarziNaftchali2018}. It reproduced field soil-water content in Canadian studies (accounting for 67--76\% of the variance) and clarified drainage mechanics, showing for compacted histosols that matric-potential drawdown was on average 58\% faster below than above a compact layer, so that raising saturated conductivity of that layer (by up to a factor of 50) rather than narrowing spacing or deepening drains is the effective remedy \parencite{Mante2018,Guedessou2021}. HYDRUS-2D simulated higher drainage rates than the Hooghoudt and Murashima--Ogino equations for paddies and highlighted the strong control of drain spacing, drain depth and cracks on drainage rate and midpoint head \parencite{Ebrahimian2015}. Related variably saturated codes include the CSUID 3D finite-difference model, which reproduced 1D infiltration, 2D water-table and 2D drainage benchmarks ($R^2=0.96$--0.99) and a conservative-transport analytical solution ($R^2=0.95$) while quantifying how parameter uncertainty propagates to root-zone salinity and crop yield \parencite{Alzraiee2013b,Alzraiee2013}, and Richards- or Darcy-based simulations of sheetpipe systems showing strong textural control of drawdown and cumulative seepage \parencite{Saptomo2021,Li2021}.

\subsubsection{Comparative and Multi-Model Evaluations}

Head-to-head comparisons generally find several field models perform comparably well. In the Netherlands, WAVE, SWAP and \textsc{DRAINMOD} predicted lateral drainage flux equally reliably, with \textsc{DRAINMOD} giving the most accurate discharge estimate \parencite{ElSadek2001}. \textsc{DRAINMOD} and SWACROP both predicted midspan water tables and drain outflow close to measurements and to each other in Atlantic Canada without calibration, \textsc{DRAINMOD} being slightly better, though both over-predicted outflow after dry spells \parencite{Prasher1996}. The SWAP model reproduced water-table behaviour and cumulative drain discharge in India with small absolute deviations (0.17--0.18 m water table; 7.8--39.4 mm cumulative discharge) and was endorsed as a simple hydraulic-performance tool \parencite{Verma2014}. \textsc{WaSim} predicted water-table depth well ($R^2=0.86$, NSE $=0.98$) but discharge less confidently ($R^2=0.57$) when driven by Rosetta-estimated conductivity, which compromised its use as a design tool \parencite{Malota2022}; independent Indian and South African applications nonetheless judged it a workable design and reclamation tool \parencite{Hirekhan2007,Malota2024}. Comparisons among distributed models show DRAINMOD tending to over-predict and CATHY to under-predict tile flow, both performing poorly at low observed flow but well during wet periods, with the soil layer below the drains identified as the most influential \parencite{Muma2017}. For salinity, a saturated--unsaturated ``SD'' approach outperformed \textsc{DRAINMOD-S} on water-table, drainage-salinity and groundwater-salinity dynamics, the latter systematically over-estimating short-term and under-estimating long-term salinity \parencite{Nozari2021}; \textsc{DRAINMOD-S} itself calibrated satisfactorily for water table, discharge and salinity in coastal China \parencite{Feng2018}.

\subsubsection{Erosion, Solute and Specialised Field Models}

Broader agricultural models have been adapted to tile drainage with mixed but generally acceptable results. \textsc{RZWQM}/\textsc{RZWQM2} simulated weekly drain flow moderately well (NSE 0.57--0.64) and reproduced atrazine transport to drains under different tillage systems, with macropore number and atrazine partitioning the key inputs \parencite{Malone2014,Singh1996}; a modified \textsc{RZWQM2} that let redistribution and tile drainage proceed simultaneously with rainfall improved daily drainage at Iowa and Ontario sites, matching cumulative drainage to within about 6--7\% \parencite{Xian2017}. \textsc{WEPP} enhancements improved uncalibrated drainage (daily NSE 0.27--0.47, rising to 0.59--0.73 after calibration) and reproduced storm and peak runoff before and after drainage installation \parencite{RevueltaAcosta2021,SAVABI1993}. \textsc{ADAPT} calibration reduced two-year cumulative-drainage under-prediction from 14\% to 7\% and performed reasonably in cold Minnesota conditions \parencite{Mendez2004}. \textsc{GPFARM-Range} simulated annual drainage, soil-water storage and content acceptably (drainage NSE $=0.76$) in subsurface-drained pastureland \parencite{Chen2016}. The \textsc{SIDRA} family, based on few parameters, has over three decades demonstrated predictive capacity for peak flows, midspan water tables and, when coupled (\textsc{SIDRA-RU}, \textsc{SIDRA-MAGE}), for catchment routing and solute transfer across French sites, with KGE $>0.81$ (calibration) and drain-flow-start median error of 5 days \parencite{Zimmer2023,Henine2022,LESAFFRE1988,Giraud1997}. Data-acquisition and diagnostic studies underpin these evaluations: automated logger--flowmeter systems measured drain flow to within 2\% of manual readings while resolving rapid event dynamics \parencite{Kanwar1999}, and combined pressure-plus-discharge measurement was shown to diagnose functional decline (a tenfold drop in potential vertical flow) in aged systems \parencite{Nishida2020,Gurklys2008}.

\subsection{Saturated versus Unsaturated Flow and Vertical Flow Resistance}

A central mechanistic finding is that treating drainage as purely saturated flow overstates discharge. A saturated flow-based model over-predicted seepage and drainage time by roughly an order of magnitude relative to experiment, an analytical solution and the SUTRA Richards-equation model, which agreed to within 20\% (seepage) and 40\% (drainage time); a vadose zone retaining substantial water above the phreatic surface persisted even at steady state, and performance depended strongly on capillary-fringe thickness, capillary rise, saturated conductivity and drain elevation \parencite{Xin2016b}. Consistent with this, drainage under saturated--unsaturated conditions is governed by gravitational, pressure, hysteresis, capillary-barrier and air-pressure potentials, so saturated conductivity alone is an inadequate design basis; water-holding capacity, permeability and aeration jointly control efficiency, and preferential drainage can become discontinuous and inefficient \parencite{Li2021}. Field diagnosis in Japanese rotational paddies showed that vertical flow resistance, not horizontal resistance, limits the maximum drainage rate: potential vertical flow (2.1--5.7 mm h$^{-1}$) was far below potential horizontal flow (6.1 to over 79 mm h$^{-1}$) \parencite{Nishida2020}. Accounting for vertical head loss through a modified procedure improved water-table prediction in paddy fields, differing from field profiles by 9--13\% versus 23--75\% for the conventional 2D solution \parencite{Talukolaee2022}. Neglecting unsaturated-zone buffering also severely underestimates response times: an analytical model that retained it matched a numerical model closely (steady heads within $0.6\pm2.0$ cm; real-case water levels NSE $=0.80$, discharge NSE $=0.89$), whereas ignoring the buffer distorted transients \parencite{VandeCraats2021}. Equivalent-medium and layered representations reproduce integrated outflow and mean water level well but can misrepresent internal dynamics, being least satisfactory where high unsaturated-zone conductivity permits fast vertical flow to the drains \parencite{Carlier2007}.

\subsection{Analytical and Semi-Analytical Solutions for Drain Spacing and Water-Table Dynamics}

\subsubsection{Steady-State versus Unsteady-State Spacing}

A robust result is that unsteady-state (transient) drain-spacing equations yield markedly narrower spacings than steady-state equations, and are preferred where extreme short-duration water-table depths occur in high-storage systems \parencite{OOSTERBAAN1988}. In Iranian paddies, average computed spacing was 30 m for unsteady-state versus 55 m for steady-state equations (26 m Glover--Dumm to 63 m Ernst), the unsteady solutions being judged more suitable given highly variable rainfall \parencite{DarziNaftchally2014}. Steady Hooghoudt-based estimates did not match observed spacings well in several settings \parencite{Moshayedi2020}, and among transient options the Glover--Dumm and van Schilfgaarde models were optimal for different rainfall durations \parencite{Moshayedi2020}. Design-criterion choice matters: for a 25 cm drop in 2 days a newly developed equation estimated spacing closest to the actual 50 m (3.76\% deviation), whereas for longer drainage periods the Glover equation was superior \parencite{Pali2014}. Nondimensional analyses show discharge following a common S-shaped curve against drain-spacing ratio, with the median at a spacing-to-depth ratio of 2--5 \parencite{Nédélec2010}.

\subsubsection{Evapotranspiration, Recharge and Sloping or Layered Aquifers}

Incorporating evapotranspiration and unsaturated-zone flow generally permits wider, more economical spacing. Analytical solutions accounting for ET lower the water table faster and allow spacing increases of about 9--18\% (constant $K=3$ m day$^{-1}$, ET $=8$ mm day$^{-1}$), applicable even to layered profiles \parencite{NIKAM1992,Singh2010,Upadhyaya2000}. Accounting for evaporation from bare soil widened calculated spacings by 23--74\% depending on texture (e.g. clay from 22.5 to 28 m; silty clay loam from 23 to 40 m) \parencite{AbdelFattah2005}. A spacing equation including horizontal unsaturated-zone flow above the water table gave up to 30\% wider spacing for coarse soils near an impervious layer, diminishing to under 5\% when that layer lay 3 m below drain level, and agreed with Hydrus-2D for non-submerged drains \parencite{Yousfi2014}. Bi-level drainage solutions establish that the deep drain discharges more than the shallow one, that including ET increases spacing, and that spacing is only slightly smaller than an equivalent conventional level system \parencite{Verma1998,Upadhyaya2000,Kankarej2023}. For sloping aquifers, water-table profiles become asymmetric with the divide shifting toward the lower drain, and hybrid finite-analytic solutions yield 12--17\% larger spacing than analytical solutions \parencite{Singh2010,Upadhyaya2013,Hartani2001}. Dual-depth systems always remove water more effectively than single-depth systems, with relative discharge reaching roughly 160--270\% depending on the depth of added tubes, so retrofitting intermediate drains can raise effectiveness by 50\% or more \parencite{Kirkham1997}.

\subsubsection{Layered Soils, Ponded Ditch Drainage and Benchmark Verification}

Analytical solutions for multi-layered and ponded-ditch systems have been verified against experiments, prior analytical work and MODFLOW with close agreement, establishing that ponded drainage is strongly three-dimensional near the drains and highly sensitive to directional conductivities and anisotropy \parencite{Chakrabarti2020,Sarmah2015,SHARMA1991,Afruzi2014}. A plow-sole layer greatly restricts seepage and lengthens particle travel times, which can be reduced by staggered ponding distributions and lower ditch levels \parencite{Chakrabarti2020,Sarmah2015}. The conductivity of even a thin lower layer significantly affects spacing \parencite{SHARMA1991}. Time-stepping linearisation of the Boussinesq equation was reported to outperform earlier Werner and Baumann techniques \parencite{Lin2025}, and iteration schemes for non-homogeneous soils reproduced field water-table responses adequately \parencite{CHU1977}. On slopes up to about 0.15 m m$^{-1}$, flat-land drainage theory remains applicable for interior parallel drains, with deviations appearing only at steeper slopes \parencite{FIPPS1989}. Cross-model verifications (e.g. SIDRA versus SLOP; analytical versus finite-element) generally show water-table shape-factor differences below 4\% for flat lands, with numerical instability arising for steep slopes \parencite{LORRE1994,Nédélec2010}. Critiques note, however, that a single linear-diffusion model cannot describe all stages of transient drawdown, since flow is radial near the tubes at early time and quasi-horizontal at late time \parencite{Kacimov2000}.

\subsection{Model Structure and Preferential-Flow Representation}

At field to hillslope scale, the representation of macropores and preferential flow is decisive for reproducing rapid drainage. Dual-permeability and dual-porosity models reproduced drainage hydrographs well (calibration $R^2=0.84$), the dual-permeability variant giving the best groundwater-level match ($r=0.98$) and internal flow representation, while a single-porosity model achieved acceptable accuracy with far fewer parameters; all approaches failed to match low continuous dry-season discharge \parencite{Varvaris2018}. Dual-permeability formulations enable use of low matrix conductivities while still reproducing fast measured drainage events through macropore flow, and are well suited to complex topography and structured soils \parencite{Warsta2014,Salo2017}. In clay-dominated hillslopes, evapotranspiration and macropore flow to the drains dominate the water balance while Darcian matrix flow is negligible, though models can struggle to dry the profile in late summer \parencite{Koivusalo1999}. Adding macropore flow in a MIKE SHE application did not change total drainage but improved peak-flow simulation, with the drainage time constant (best value 5.6 h) the most critical parameter \parencite{Zhou2013}. Fully coupled surface--subsurface--tile models are recommended for field-scale dynamics, with discrete pipes most relevant for drain-flow dynamics and dual-permeability meshes reserved for detailed contaminant-transport studies \parencite{DeSchepper2015}. Representing the drain network as an equivalent porous medium or via seepage nodes can reproduce outflow almost exactly even when internal heads differ, offering a practical simplification for watershed integration \parencite{DeSchepper2015,Carlier2007,Boico2022}. Simply demonstrated is that letting drainage occur simultaneously with rainfall and redistribution improves daily and peak-timing skill \parencite{Xian2017}, and that CFD analysis of corrugated-pipe geometry (outside diameter the dominant control) matched sand-tank experiments to within 6\% \parencite{Wu2022}.

\subsection{Watershed- and Regional-Scale Drainage Representation}

\subsubsection{Sensitivity of Streamflow Partitioning to Drain Conceptualisation}

How drains and subsurface heterogeneity are represented strongly controls simulated catchment behaviour. Two grid conceptualisations of the same Danish catchment produced drain contributions to streamflow of 13\% (regular grid) versus 31\% (irregular grid), with the irregular model better capturing the lowest measured groundwater--surface-water exchange fluxes despite both matching total streamflow within 0.1\% \parencite{Karan2021}. Comparing conceptual models built on seepage nodes, soil heterogeneity affected drainage-area simulations more than tile-drain conceptualisation, while the reverse held at subcatchment scale and both mattered comparably at catchment scale; seepage-node models simulated similar discharges and, being about 35\% faster, are better candidates for automated calibration and for cases where drain locations are unknown \parencite{Boico2022}. Distributed HydroGeoSphere models reproduced observed stream hydrographs and flashy heads better than MIKE SHE, which produced slowly varying heads and lower peak frequencies \parencite{DeSchepper2017}. Machine-learning work established that streamflow signatures, especially high-flow event duration, predict simulated drain fraction across 73 Danish sub-catchments (test MAE $=0.126$), and that including signatures in calibration improved fits without major trade-offs \parencite{Schneider2022}. Novel strategies that embed artificial drainage in fully distributed integrated hydrological models reproduced streamflow without site-specific calibration, outperforming a calibrated SWAT model \parencite{Rathore2024}. At watershed outlets, 3D CATHY simulations showed subsurface drainage raising cumulative outlet flow (by 3--22\% across scenarios), lowering peak and raising minimum flows, with exfiltration the dominant outlet component \parencite{Muma2016}. Watershed-scale water balances also quantify drain contributions: subsurface drainage was the dominant source (about 68\% of total drainage) in a Chinese plain irrigation district \parencite{Zheng2025b}, and drainage represented major outflow fractions of precipitation across Nordic monitored catchments \parencite{Nousiainen2015,Turunen2013}.

\subsubsection{MODFLOW-Based Regional Drainage Design}

Distributed groundwater models, principally MODFLOW, are established tools for regional drainage assessment, salinity-hazard mapping and design optimisation. Calibrated MODFLOW models matched observed heads closely (e.g. 92\% correlation; standard errors of order 0.1--0.2 m) and were used to design interceptor and additional-drain systems, quantify drawdown (average 0.53 m over three years) and predict discharge for candidate designs \parencite{Ghorbani2016,Mirlas2012,Mirlas2009,Mirlas2013,Mirlas2021,AbdElaziz2025}. Sensitivity analyses repeatedly showed these regional models to be relatively insensitive to specific storage ($\pm15\%$) and aquifer conductivity ($\pm50\%$) but responsive to the conductivity of aquitard/relative-aquiclude layers \parencite{Mirlas2012,Mirlas2009,Mirlas2013}. Water-balance analyses attributed distinct fractions of flow to drainage, evaporation and seepage and revealed that existing drains were sometimes partly ineffective, intercepting as little as 9--30\% of groundwater flow \parencite{Mirlas2009,Mirlas2012}. Simulation--optimisation coupled to MODFLOW achieved equal or better water-table control with about 20\% less pumping and fewer wells \parencite{Rana2008}, and drainage-orientation studies identified a 90$^{\circ}$ alignment to groundwater flow as optimal \parencite{AbdElaziz2025}. Coastal-aquifer applications quantified how sea-level rise advances seawater intrusion and raises drain discharge by 12.5--43.4\% under future scenarios \parencite{AbdElaty2019}. New MODFLOW subsurface-drainage (SDR) packages based on the Hooghoudt equation markedly outperform the conventional DRN package: SDR reproduced steady-state and experimental drainage (RRMSE $<5\%$; relative errors mostly $<6\%$) by calibrating only conductivity, whereas fixed DRN conductance caused large errors when layout or conductivity changed, and simulated 16.9--55.7 times more drainage water than DRN \parencite{Yang2022b,Han2025}. A further modification introducing an entrance-resistance coefficient corrected large over-predictions for synthetic-envelope pipes (NRMSE $<3.18\%$; NSE and KGE $>0.9$) \parencite{Qian2025b}. Such frameworks also quantified salt dynamics, showing regional salt-discharge increases of 38--54\% and vadose-zone salt reductions exceeding 90\% after subsurface drainage, and a critical drainage-to-irrigation ratio near 10--11\% to avoid salt accumulation \parencite{Yang2022b,Han2025}. Grid resolution matters: a 1 m grid was appropriate for the SDR package, with relative errors below 4.8\% \parencite{Han2025}. Regional nitrate-reduction mapping further showed that high-drainage areas tend to have moderate-to-low groundwater reduction because drainage bypasses reducing subsurface zones \parencite{Hojberg2017,Hansen2019}.

\subsection{Sensitivity, Calibration and Uncertainty}

Formal analyses converge on saturated hydraulic conductivity and evapotranspiration estimation as the dominant sources of sensitivity and uncertainty. \textsc{DRAINMOD} was most sensitive to vertical conductivity of the restrictive layer and lateral conductivity of the deepest soil layer, with parameter interactions revealed only when likelihood weighting was applied \parencite{Wang2006}. Conductivity estimation method produced ``tangible'' differences in simulated water-table fluctuations, the direct method giving best agreement \parencite{NabiSichani2012}, and the van Genuchten shape parameter $n$, lateral conductivity and depth to the impermeable layer ranked as the leading sensitive parameters in HYDRUS-based inverse modelling \parencite{Sedaghatdoost2018}. Inverse modelling with a genetic algorithm robustly estimated hydraulic parameters, and monitoring drainage discharge together with water-table depth was found sufficient (NSE up to 0.83--0.95) \parencite{Sedaghatdoost2018}. Uncertainty frameworks combining GLUE, extended FAST and \textsc{DRAINMOD} showed annual predictions falling within confidence bounds and exposed interaction effects invisible to one-at-a-time analysis \parencite{Wang2006}. Multi-objective sensitivity analysis of 17 hydrologic parameters found daily drainage flow and relative corn yield most sensitive to drain spacing and drain depth, with soil-temperature parameters affecting only winter water balance, underscoring the need to calibrate the most sensitive parameters carefully \parencite{Timalsina2023}. For drainage-water-recycling design, the soil-water control volume was the most influential input, guiding interface simplification \parencite{Reinhart2020}. Parameter uncertainty was shown to propagate materially into predicted root-zone salinity, water-table depth and relative crop yield, motivating probabilistic rather than deterministic design \parencite{Alzraiee2013,Alzraiee2013b}. Consistent with this, the review literature concludes that no universal model exists, that automatic calibration is preferable because manual calibration neither guarantees global optima nor supports uncertainty estimation, and that integrated multi-model use is often more practical than single-model use \parencite{Singh2019b}.

\subsection{Design Applications, Field Outcomes and Supporting Measurements}

The validated models have been applied extensively to quantify design trade-offs. Narrower spacing and shallower drains improve water-table control and aeration: SEW30 fell sharply from 3461 to 450 cm day across configurations as drains were made shallower and closer in Iranian paddies, and a 0.65 m depth with 30 m spacing provided desired wet-season control \parencite{DarziNaftchali2013,DarziNaftchally2014}. Design analyses balanced aeration against drainage volume, showing SEW30 rising and drainage volume falling with wider spacing, and mild-to-severe aeration risk at 75 m but not at 25 m spacing \parencite{Verma2014,Hirekhan2007}. Optimum spacings for trafficability and crop protection on Illinois claypan soils were 10--13 m \parencite{MOSTAGHIMI1989}, and in South Africa target mean water tables of 1.0--1.5 m required 25--40 m spacing in clay versus 55--70 m in clay loam, confirming the closer-spacing-for-lower-$K_{sat}$ design rule \parencite{Malota2015}. Shallow and controlled drainage were shown to reduce drainage volume: shallow drainage cut Iowa subsurface drainage by 8.4\% \parencite{Zhou2013,Mendez2004}, and controlled drainage starting 1 March raised groundwater tables by 19--35 cm and reduced annual outflow by 35--52 mm and nitrate losses by 15--22 kg ha$^{-1}$ across climate years \parencite{Kesicka2023}, while lateral seepage explained 96\% of the drainage reduction under controlled drainage in North Carolina, so that ignoring seepage severely over-predicts outflow \parencite{Youssef2021,Saadat2018,Österholm2015}. Climate-change simulations consistently project higher future subsurface drainage (increases of about 14--80\% depending on model and precipitation change) concentrated in winter, with implications for nitrate loss, while recommended designs still maintained relative yields above 80\% \parencite{Singh2009,Chen2016,Feng2018}. Supporting field studies established the reliability of measurement and property-estimation methods underpinning all this modelling: drain-derived hydraulic conductivity and drainable porosity agreed with independent methods \parencite{BUCKLAND1986}, Rosetta and RETC reproduced measured conductivity and retention accurately \parencite{Malota2022,Malota2024}, cover crops reduce drainage by a robust $\approx30$ mm \parencite{Meyer2019}, and long-term drainage intensity measurably alters soil physical properties, carbon accumulation and workability, with $\le12$ m spacing protecting subsoil from compaction \parencite{Omonode2025,Mante2022,Mante2018,Frison2009,Roy2020}. Controlled-drainage research has grown rapidly, and bibliometric analysis of 462 articles (2000--2022) confirms its emergence as a leading, multidisciplinary management practice for water quality in tile-drained landscapes \parencite{Kesicka2023b}. Emerging model-predictive control and drainage-water-recycling tools extend these capabilities toward real-time water-table management and anticipatory drainage-base control \parencite{Ikonen2023,Reinhart2020,VandeCraats2021,Ahmadi1995,HASEGAWA1994,GALLICHAND1992b,MCCARTHY1992,PalaciosVélez2004,Khan2007,Prathapar1996,Ghorbani2015,Skaggs2017,Singh2006,Singh1999,Gaj2022}.